\chapter{Interrupts}
\label{ch:interrupts}

The context switch in the previous chapter does not run by magic; something
has to \emph{call} the scheduler. On this port that something is almost always an
interrupt: the systimer tick, or a device interrupt whose handler makes a
higher-priority task ready. This chapter follows an interrupt from the silicon
vector all the way to an Ada handler, and explains how \texttt{pragma
Attach\_Handler} is supported across all three profiles, including the
bare-board full \texttt{System.Interrupts} the \texttt{full} profile needs,
which no stock bare-board runtime provides.

\section{The Xtensa interrupt levels}\index{interrupts}\index{interrupt levels}\index{xt\_highint5}

The ESP32-S3's Xtensa LX7 has six maskable interrupt levels plus a non-maskable
one, and two of those levels are special:

\begin{longtable}{p{2.2cm} p{10.5cm}}
\toprule
\textbf{Level} & \textbf{Role on this port} \\ \midrule
\endhead
1     & The general \emph{exception} vector (synchronous faults). No async
        interrupt is assigned here. \\
2--3  & \textbf{The OS band} (\texttt{\_\_gnat\_level2\_vector} /
        \texttt{\_\_gnat\_level3\_vector}): every interrupt that talks to
        the scheduler. The device interrupts live here, and so do the systimer
        tick and the cross-core IPI ``poke'' (both share CPU interrupt 21, at
        level 2). \\
4--5  & \emph{Free for OS-oblivious handlers}: lean assembly that never
        touches the scheduler. Nothing is installed today; the old level-5
        vector (\texttt{xt\_highint5}) survives only as a park stub. \\
6     & \texttt{DEBUGLEVEL} (debug exceptions: breakpoints and the
        \texttt{DBREAK} data watchpoints, used for stack-overflow detection). \\
7     & \texttt{NMILEVEL} (non-maskable; not used). \\
\bottomrule
\end{longtable}

One architectural constant shapes everything below.
\texttt{XCHAL\_EXCM\_LEVEL} is 3, and it draws the line the whole design obeys:
\emph{every interrupt that saves context or can trigger a task switch sits at
or below it}. \texttt{PS.EXCM} -- which the hardware sets during every window
over/underflow exception -- masks exactly levels 1--3. So does the raised
\texttt{INTLEVEL} across every register-window spill. Keep all the
scheduler-facing interrupts inside that band and no interrupt can ever observe
a half-spilled register window. This is the FreeRTOS Xtensa port's rule, and
this port adopted it deliberately (verified against the esp-idf source). The
original design ran the tick at level 5, above the line -- and a level-5 tick
can land in the middle of window-critical code and corrupt the register state.
Levels above \texttt{EXCM\_LEVEL} remain available, but only for handlers
that never touch the scheduler.

\section{Priorities and the masking model}\index{interrupt priorities}\index{interrupt masking}

Ada interrupt priorities map onto Xtensa levels in
\texttt{System.BB.Board\_Support.Priority\_Of\_Interrupt}: a hardware interrupt's
fixed level becomes the Ada priority its handler runs at, and conversely a task's
ceiling becomes the \texttt{INTLEVEL} it raises while it holds a lock
(\texttt{Enable\_Interrupts} in \texttt{s-bbcppr}). \texttt{Ada.Interrupts.Names}
gives symbolic IDs for the slots that are free for application use:
\texttt{Device\_L2\_0/1} (CPU interrupts 19/20) and \texttt{Device\_L3\_0/1}
(23/27), with ceilings \texttt{Device\_L2\_Priority} /
\texttt{Device\_L3\_Priority}. CPU interrupt 21 (\texttt{Device\_L2\_2}) is
kernel-reserved: the systimer tick and the cross-core poke share it, and the
kernel attaches its alarm handler there at start-up. Do not attach to it.

\section{What the silicon does}\index{EPC}\index{EPS}\index{EXCSAVE}\index{VECBASE}\index{rfi}

Before any of our code runs, the CPU itself performs a fixed sequence. It is
small, and everything above depends on it, so here it is in full. When an
enabled interrupt of level $n$ (2--7) is the highest one pending:

\begin{enumerate}
  \item The hardware copies the current \texttt{PC} into \texttt{EPC$n$} and
        the current \texttt{PS} into \texttt{EPS$n$}. Each level has its own
        pair, so a level-3 interrupt cannot clobber the state a level-2 entry
        is still using.
  \item It sets \texttt{PS.INTLEVEL} to $n$, masking level $n$ and below.
  \item It jumps to the level's fixed slot in the vector table:
        \texttt{VECBASE} plus a per-level offset. Nothing is pushed on any
        stack --- there is no hardware stacking on Xtensa. The vector gets the
        CPU exactly as the interrupted code left it, registers and all.
\end{enumerate}

That last point drives the vector's first moves. The slot has \emph{no} free
register, so its first instruction stashes \texttt{a0} in the level's
\texttt{EXCSAVE$n$} scratch register; only then can it start saving state
properly. From C this is the part the NVIC did for you on a Cortex-M ---
here the eight-instruction prologue that builds the frame is ours, in
\texttt{highint5.S}.

The return trip is one instruction: \texttt{rfi $n$} restores
\texttt{PS} from \texttt{EPS$n$} and jumps to \texttt{EPC$n$}
\emph{atomically}. That atomicity is load-bearing: \texttt{INTLEVEL} stays
at $n$ until the very last instruction, so there is no window in which a
fresh interrupt could preempt a half-restored context.

Two registers complete the picture. \texttt{INTERRUPT} (read-only) shows
which of the 32 CPU interrupts are pending right now; the dispatch reads it
to find the sources to service. \texttt{INTENABLE} gates them per core; the
runtime sets a bit when a handler is attached. Most sources are
\emph{level-triggered}: the CPU-side bit stays asserted until the handler
clears the \emph{peripheral's} status register (usually write-1-to-clear),
which is why every HAL handler acks the device, not the CPU.

\section{The interrupt matrix}\index{interrupt matrix}\index{INT\_MAP}

If you come from a Cortex-M, one ESP32 idea needs installing first: peripheral
interrupts do not have fixed CPU positions. The chip has roughly a hundred
peripheral \emph{sources} (TWAI, each GDMA channel, the systimer comparators,
GPIO, \dots) and each core has just 32 CPU interrupts, each with a
\emph{fixed, hardwired} level and trigger type. Between the two sits the
\emph{interrupt matrix}: one map register per source, per core
(\texttt{INTERRUPT\_CORE0\_*\_INT\_MAP}), holding the number of the CPU
interrupt that source should raise.

Three consequences matter in practice:

\begin{itemize}
  \item \textbf{Routing is a design decision.} A driver picks its
        interrupt's level by picking which CPU interrupt to map its source
        to. That is how the LCD refill moved from level 2 to level 3, and how
        the tick moved from level 5 to level 2 --- one store to a map
        register each.
  \item \textbf{Sources can share.} Several sources may map to one CPU
        interrupt; the handler then reads the peripherals to tell them apart.
        The kernel uses this deliberately: the systimer alarm and the
        cross-core poke share CPU interrupt 21, and the dispatch checks the
        \texttt{FROM\_CPU} register to distinguish them.
  \item \textbf{Not every CPU interrupt is usable.} The 32 slots are a mixed
        bag: some are level-triggered, some edge-triggered, some are internal
        (software) interrupts the matrix cannot drive at all. The usable
        level-triggered OS-band slots are few --- 19/20/21 at level 2, 23/27
        at level 3 --- which is why the allocation in
        \texttt{Ada.Interrupts.Names} is as tight as it is. (Two findings
        from the field: CPU interrupt 22 is edge-triggered and will simply
        never fire for a level source, and 29 is an internal software
        interrupt the matrix cannot reach. Both look free on paper.)
\end{itemize}

\section{From silicon to handler}\index{interrupt vector}\index{interrupt wrapper}

\texttt{start.S} points \texttt{VECBASE} at our vendored vector table before
anything else runs, so every vector is ours. The level-2 and level-3 vectors
are two instantiations of one assembly macro (\texttt{gnat\_os\_vector} in
\texttt{highint5.S}), and they mirror the FreeRTOS Xtensa port's
\texttt{\_frxt\_int\_enter}/\texttt{\_frxt\_int\_exit} structure. When an
OS-band interrupt fires (say the tick, on CPU interrupt 21), this is the path:

\begin{enumerate}
  \item \textbf{Vector entry.} The CPU jumps to the level-2 slot. The
        vector's first instruction masks to \texttt{DEBUGLEVEL}
        (\texttt{rsil 6}): the whole prologue is now atomic against the OS
        band, and the frame build below cannot trip the per-task
        stack-overflow watchpoint.
  \item \textbf{Save.} It builds an \texttt{XT\_STK} frame on the
        interrupted stack and calls \texttt{\_xt\_context\_save} to spill the
        full register state into it. The spill leaves
        \texttt{WINDOWSTART = 1 << WINDOWBASE} by construction, so a plain
        resume needs no window bookkeeping at all.
  \item \textbf{Enter.} Still masked, it increments the per-CPU nesting
        count (\texttt{\_\_gnat\_int\_nest}, owned entirely by the vector).
        On the outermost entry it saves the task-side SP and hops onto the
        per-CPU \emph{interrupt stack}, so handler frames and their spills
        never land on task stacks.
  \item \textbf{Service.} It sets \texttt{INTLEVEL} to
        \texttt{XCHAL\_EXCM\_LEVEL} and calls the Ada dispatch routine
        (\texttt{\_\_gnat\_level2\_dispatch} or
        \texttt{\_\_gnat\_level3\_dispatch}). The dispatch reads the Xtensa
        \texttt{INTERRUPT} register and, for each pending source it owns,
        calls \texttt{System.BB.Interrupts.Interrupt\_Wrapper (Id)}. Because
        every handler runs at \texttt{EXCM\_LEVEL}, OS-band handlers never
        preempt one another; the interrupt \emph{level} decides which pending
        source is taken first and which interrupts preempt task code.
  \item \textbf{The wrapper} (\texttt{s-bbinte.adb}) is the GNARL bridge: it
        records the interrupt being handled, raises the running task's active
        priority to the interrupt's priority (so the ceiling protocol holds),
        runs the \emph{registered handler}, then restores priority. The
        handler is an \texttt{access procedure (Id : Interrupt\_ID)}.
  \item \textbf{Epilogue: the one switch site.} Back in the vector, masked
        again, the nesting count is decremented. Only the \emph{outermost}
        exit may switch, and it decides by itself: it compares
        \texttt{First\_Thread} with \texttt{Running\_Thread} and, if they
        differ, tail-jumps \texttt{\_\_gnat\_preempt\_dispatch} from the
        vector's clean single-window context. One guard applies: if the
        interrupted PC lies inside vector text, the switch is deferred to the
        next interrupt exit -- a half-entered vector's \texttt{EPC/EPS/%
EXCSAVE} registers are still live in hardware and must not be stranded.
  \item \textbf{Resume.} Otherwise \texttt{\_xt\_context\_restore} +
        \texttt{rfi} returns to the interruptee atomically:
        \texttt{INTLEVEL} stays high until the very last instruction, so
        there is no re-entrancy window.
\end{enumerate}

The level-3 vector is the same macro with \texttt{EPS\_3/EPC\_3} and
\texttt{rfi 3}. Note what is \emph{not} here: no per-vector deferral flags,
no software \texttt{WINDOWSTART} repair on the plain-resume path, and no
special tick vector -- the tick is just another level-2 source, dispatched
like any device interrupt.

\section{Attaching a handler: two axes}\index{Attach\_Handler}\index{protected handler}

It helps to keep two \emph{independent} distinctions apart. The first is
\emph{which API layer} you attach through; the second (which lives entirely
inside the high-level layer) is \emph{static versus dynamic} attachment.

\subsection{Axis 1: API layer (high level vs low level)}\index{Ada.Interrupts}

\textbf{High level: \texttt{Ada.Interrupts} / \texttt{pragma Attach\_Handler}.}
The portable Ada way: the handler is a protected procedure, the enclosing
protected object's \texttt{Interrupt\_Priority} is its ceiling, and the compiler
generates the registration and ceiling locking. This is what the GPIO HAL and the
\texttt{esp32s3\_intr\_levels} example use. The same counting handler as the
low-level version below, written this way, is:

\begin{lstlisting}
with Ada.Interrupts.Names; use Ada.Interrupts.Names;

package body High is
   protected PO with Interrupt_Priority => Device_L2_Priority is
      function Count return Natural;
   private
      procedure Handle;
      pragma Attach_Handler (Handle, Device_L2_0);  -- installed at elaboration
      N : Natural := 0;
   end PO;

   protected body PO is
      procedure Handle is           -- runs at the PO ceiling; locking is implicit
      begin
         Clear_Source;              -- ack the (level-triggered) source
         N := N + 1;
      end Handle;
      function Count return Natural is (N);
   end PO;
end High;
\end{lstlisting}

There is no \texttt{Install} call: \texttt{pragma Attach\_Handler} wires the
handler up when \texttt{PO} elaborates. The handler is a protected procedure (no
\texttt{Id} parameter, since it is bound to one interrupt), the count \texttt{N} is a
protected component rather than a bare \texttt{Volatile}, and the ceiling protocol
serialises it against other accesses to \texttt{PO} for free. Contrast the
hand-rolled low-level form:

\textbf{Low level: \texttt{System.BB.Interrupts.Attach\_Handler}.} The
bareboard kernel's own primitive: register a plain \texttt{access procedure (Id)}
for an \texttt{Interrupt\_ID}, with an explicit priority. This is what
\texttt{Interrupt\_Wrapper} ultimately calls, it is available on \emph{all}
profiles, and it needs no protected object, but also gives you no
compiler-enforced ceiling. It is the escape hatch when the portable layer is
unavailable or unwanted:

\begin{lstlisting}
with System; use System;
with System.BB.Interrupts; use System.BB.Interrupts;

package body Low is
   Count : Natural := 0 with Volatile;

   procedure Handle (Id : Interrupt_ID) is
   begin
      Clear_Source (Id);          --  ack the (level-triggered) source
      Count := Count + 1;
   end Handle;

   procedure Install is
   begin
      --  plain procedure, explicit Id + priority, no protected object
      Attach_Handler (Handle'Access, Id => 19,
                      Prio => Interrupt_Priority'Last);
   end Install;
end Low;
\end{lstlisting}

\subsection{Axis 2: static vs dynamic (within the high-level layer)}\index{dynamic attachment}

Static versus dynamic is the dividing line of the \texttt{No\_Dynamic\_Attachment}
restriction; it only concerns the high-level \texttt{Ada.Interrupts} layer.

\textbf{Static: \texttt{pragma Attach\_Handler}.} The handler is fixed at
compile time and installed once, automatically, when the protected object is
elaborated. There is no run-time attach call; the compiler lowers it straight to
\texttt{System.Interrupts.Install\_Handlers}, \emph{bypassing}
\texttt{Ada.Interrupts}. This is the form to reach for by default:

\begin{lstlisting}
with Ada.Interrupts.Names; use Ada.Interrupts.Names;

package body Blink is
   protected L2_PO with Interrupt_Priority => Device_L2_Priority is
   private
      procedure Handle;
      pragma Attach_Handler (Handle, Device_L2_0);  --  installed at elaboration
      N : Integer := 0;
   end L2_PO;

   protected body L2_PO is
      procedure Handle is
      begin
         Clear_L2;                 --  ack the source
         N := N + 1;
      end Handle;
   end L2_PO;
end Blink;
\end{lstlisting}

\textbf{Dynamic: the \texttt{Ada.Interrupts} operations.} Here the handler is
attached, swapped, or detached \emph{at run time}. The protected procedure must
first be made eligible with \texttt{pragma Interrupt\_Handler}; then
\texttt{Attach\_Handler}, \texttt{Detach\_Handler}, \texttt{Exchange\_Handler} and
\texttt{Current\_Handler} act on it during execution. These calls go \emph{through}
the \texttt{Ada.Interrupts} body: the path the bare runtime initially stubbed
out (every operation raised \texttt{Program\_Error}) until a delegating
\texttt{a-interr} body was supplied (see \emph{Why the profiles differ} below):

\begin{lstlisting}
with Ada.Interrupts;       use Ada.Interrupts;
with Ada.Interrupts.Names; use Ada.Interrupts.Names;

package body Dyn is
   protected PO is
      procedure Handle;
      pragma Interrupt_Handler (Handle);   --  eligible for dynamic attach
   end PO;
   protected body PO is
      procedure Handle is begin null; end Handle;
   end PO;

   procedure Run is
   begin
      Attach_Handler (PO.Handle'Access, Device_L2_0);  --  attach now
      --  ... later, restore the default treatment ...
      Detach_Handler (Device_L2_0);
   end Run;
end Dyn;
\end{lstlisting}

So the two axes compose: a high-level handler is either static (the \texttt{pragma},
installed at elaboration) or dynamic (the run-time \texttt{Ada.Interrupts} calls),
while the low-level \texttt{System.BB.Interrupts} primitive sits beneath both,
itself a run-time call, but \emph{not} the Ada \texttt{No\_Dynamic\_Attachment}
notion of ``dynamic''.

\section{Why the profiles differ}\index{runtime profile}\index{System.Interrupts}

The high-level form compiles differently depending on the tasking profile, and
that is the crux:

\begin{itemize}
  \item \texttt{embedded} and \texttt{light-tasking} set \texttt{pragma Profile
        (Jorvik)}. Under a Ravenscar-family profile the compiler lowers
        \texttt{pragma Attach\_Handler} to the \emph{restricted} installation path,
        \texttt{System.Interrupts.Install\_Restricted\_Handlers}, which the
        bareboard \texttt{s-interr} provides. So it works.
  \item \texttt{full} sets \emph{no} profile pragma (its whole point is to lift
        the Jorvik restrictions). The compiler then lowers interrupt-handler
        protected objects to the \emph{full dynamic} machinery:
        \texttt{Register\_Interrupt\_Handler}, \texttt{Install\_Handlers}, and the
        \texttt{Static\_Interrupt\_Protection} / \texttt{Dynamic\_Interrupt\_-
        Protection} protected types.
\end{itemize}

(Note this is the \emph{profile} distinction, not \texttt{Configurable\_Run\_Time},
which is \texttt{True} in both.) The catch is that AdaCore ships a full
\texttt{System.Interrupts} only for hosted (POSIX) targets, where it is built on
a signal-server task and is unportable to bare metal; \emph{no} bare-board
runtime (not even AdaCore's own \texttt{ravenscar-full}) implements the
full surface because Ravenscar forbids dynamic attachment in the first place.
For a long time, then, a \texttt{full}-profile build of an \texttt{Ada.Interrupts}
handler simply failed with \texttt{construct not allowed in this configuration}
and \texttt{System.Interrupts.Register\_Interrupt\_Handler not defined}.

\subsection{The bare-board full \texttt{System.Interrupts}}\index{umbrella handler}

That gap is now closed. \texttt{full\_overlay/gnarl/s-interr}
is a bare-board re-implementation of the full \texttt{System.Interrupts}
surface (\texttt{Register\_Interrupt\_Handler}, the \texttt{Static\_-} and
\texttt{Dynamic\_Interrupt\_Protection} protected bases, \texttt{Install\_-
Handlers} and the dynamic \texttt{Attach\_}/\texttt{Detach\_}/\texttt{Exchange\_-
Handler} operations) but layered on the machinery that already exists rather
than on a signal server. It only had to add the GNARL-facing surface; the hard
parts were already done and proven (the previous sections):

\begin{itemize}
  \item It installs a single \emph{umbrella} handler per interrupt through the
        low-level \texttt{System.OS\_Interface.Attach\_Handler}
        (= \texttt{System.BB.Interrupts}), then redirects it by re-pointing a
        handler table. The bare-board primitive accepts only \emph{one} attach
        per source, so the umbrella is installed once and the table does the rest.
  \item The interrupt's run priority is fixed by its Xtensa level, so the
        controller-programming step ignores the priority argument
        (\texttt{Install\_Interrupt\_Handler} marks it \texttt{Unreferenced}).
        That is what makes the dynamic \texttt{Attach\_Handler} path, which
        carries no priority, safe.
  \item Interrupt \emph{entries} and the POSIX.5 signal services have no
        bare-board meaning and raise \texttt{Program\_Error}.
\end{itemize}

The overlay file replaces the restricted \texttt{s-interr} for the \texttt{full}
profile \emph{only}; \texttt{embedded} and \texttt{light-tasking} keep the
restricted one they need for the Jorvik path. So the same protected-handler
source now builds on all three profiles.

There is one wrinkle, and it is subtler than it first looks. The full
\texttt{System.Interrupts} alone is necessary but not \emph{sufficient}: the
\emph{static} \texttt{pragma Attach\_Handler} path reaches it directly (the
compiler lowers the handler to \texttt{Install\_Handlers}), but the \emph{dynamic}
operations go through \texttt{Ada.Interrupts}, and the bare runtime inherited the
\emph{restricted} (\texttt{No\_Dynamic\_Attachment}) \texttt{Ada.Interrupts} body
from the Jorvik base, in which every operation is simply \texttt{raise
Program\_Error}. So for a while static attach worked while
\texttt{Ada.Interrupts.Attach\_Handler}, \texttt{Detach\_Handler},
\texttt{Exchange\_Handler}, \texttt{Current\_Handler}, \texttt{Is\_Attached} and
\texttt{Reference} all raised \texttt{Program\_Error} (the full ACATS sweep caught
this: \texttt{CXC3007}). The overlay now also supplies the \emph{delegating}
\texttt{a-interr.adb} that forwards each operation to \texttt{System.Interrupts},
so the full dynamic API works (\texttt{CXC3007} HW-verified \texttt{PASSED}).

The genuine residual limitation is narrow: because \texttt{Configurable\_Run\_Time}
is \texttt{True}, the front end does not generate library-level finalization, so a
\emph{library-level} interrupt protected object's \texttt{Finalize} (which would
restore the previously installed handlers at program shutdown) does not run
automatically. Explicit \texttt{Detach\_Handler} works; only the implicit
restore-on-finalization of a library-level handler PO is absent. That is harmless
on an embedded target that never exits, and the low-level
\texttt{System.BB.Interrupts.Attach\_Handler} primitive remains available on all
profiles as the no-finalization escape hatch.

\section{Exercising the vectors}\index{CPU\_INT}

Because a bug in the vectors (a corrupted context, a dropped interrupt) is exactly
the kind of thing that hides until the worst moment, the
\texttt{esp32s3\_intr\_levels} example is a standalone regression test for them. It
fires the level-2 and level-3 device vectors with \emph{no external wiring} by
routing the \texttt{FROM\_CPU} interrupt-matrix sources (the same mechanism as the
cross-core poke) to \texttt{Device\_L2\_0} and \texttt{Device\_L3\_0} and asserting
them from software, while a low-priority victim task holds floating-point
accumulators and a \texttt{THREADPTR} sentinel live in registers, continually
preempted by those two vectors \emph{and} the timer tick. If any vector failed
to preserve the interrupted context, an accumulator or the sentinel comes back
wrong. On hardware the handler counts climb and the corruption count stays zero,
confirming all three vectors both dispatch correctly and preserve context.
(Level~4 has no vector to test; level~1 carries no asynchronous interrupt on this
port.)

The companion \texttt{esp32s3\_full\_intr} example is the same test built under
the \texttt{full} profile. Its job is narrower: to prove that the identical
\texttt{pragma Attach\_Handler} protected objects now build, bind, link \emph{and
run} against the bare-board full \texttt{System.Interrupts} described above:
the exact source that used to fail with \texttt{Register\_Interrupt\_Handler not
defined} on \texttt{full}. On hardware it boots clean on both cores and the L2
and L3 handler counts climb together with the clean-batch counter, with zero
context-loss and zero faults. So on \texttt{full}, too, the portable
\texttt{Ada.Interrupts} protected-handler form dispatches correctly and preserves
the preempted context.
